% ===========================================================================
\clearpage
\floraon
\florathispage{layout=surround, seed=37, density=0.32, visitors=0.08}
\part{The tools}
\thispagestyle{empty}
\clearpage
\floraoff
% ===========================================================================

\section{C1: static analyzer}
\label{sec:c1}

\kst{kernscope-analyze} consumes PTX or AMDGPU assembly text and emits one
kernel report plus a suggestion list. It is a static tool: it estimates
occupancy and memory behavior from the instruction stream and the architecture
database, and it labels every estimate as a metric rather than a measurement.

\subsection{Inputs, CLI, and precedence}
\label{sec:c1-cli}

\begin{lstlisting}
kernscope-analyze <file.ptx|file.s|file.isa> [options]
  --arch <sm_XX|gfxXXXX>   required for AMDGPU; PTX defaults to .target
  --block <N>              launch block size for occupancy
  --json | --text | --html output format (default --text)
  --source <file.cu|file.hip>  source stem for the kernel label
  --version | --help
\end{lstlisting}

The analyzer recognizes \kst{sm_75}, \kst{sm_80}, \kst{sm_86}, \kst{sm_89},
\kst{sm_90}, \kst{sm_100}, \kst{sm_120}, \kst{gfx90a}, \kst{gfx942}, and
\kst{gfx1101}. A \kst{.ptx} file goes through the PTX parser; \kst{.s} and
\kst{.isa} files go through the AMDGPU parser; other extensions are
content-sniffed. A directory input may carry \kst{manifest.json} with
\kst{arch}, \kst{block_size}, and a \kst{ptxas_v_report}; command-line flags
override the manifest. Without \kst{--block}, occupancy emits the
\kst{block_size_unknown} diagnostic rather than guessing.

Register pressure prefers \kst{ptxas -v} when a report is supplied. Otherwise
the analyzer uses the declared \kst{.reg} counts plus the \kst{.local} frame,
an upper bound: ptxas can reuse registers, so the proxy can overestimate.
The AMDGPU path uses \kst{scratch_load} and \kst{scratch_store} for spills and
does not have a \kst{next_free_vgpr} value in plain assembly, so it does not
invent one.

\subsection{The ten analyses}
\label{sec:c1-checks}

\begin{table}[htbp]
\centering
\small
\begin{tabularx}{\textwidth}{@{}l X l@{}}
\toprule
Analysis & What it computes & Finding \\
\midrule
Register pressure & Per-thread registers, warp allocation, local frame. & \kst{spill}, \kst{cannot_estimate_registers} \\
Occupancy + limiter & Blocks per SM from thread, warp, smem, register, and hard limits. & \kst{occupancy}, \kst{limiter} \\
Global coalescing & 128\,B sectors, used versus transferred bytes per warp. & \kst{coalescing}, \kst{coalescing_unknown} \\
Shared banks & Bank index, distinct banks, conflict way. & \kst{bank_conflict}, \kst{bank_conflict_unknown} \\
Warp divergence & Branch target and post-dominator analysis over thread-dependent predicates. & \kst{possibly_divergent_branch}, \kst{divergent_inner_loop} \\
Barriers & \kst{bar.sync}, \kst{bar.warp}, cluster barriers, density per loop, reachability under divergence. & \kst{divergent_barrier} \\
Instruction mix & FP32/FP16/INT/LSU/SMEM/tensor/control histogram. & \kst{mix_count_*}, \kst{mix_pct_*} \\
Tensor cores & PTX \kst{mma}/\kst{wmma}/\kst{wgmma}/\kst{tcgen05}; AMD \kst{v_mfma_*}/\kst{v_wmma_*}. & \kst{no_tensor_core_usage}, \kst{tensor_ops} \\
Spills & \kst{ld.local}/\kst{st.local}, local frame bytes, AMD scratch ops. & \kst{spill}, \kst{spill_bytes} \\
Latency & Static critical-path estimate from microbenchmark-derived op latencies. & \kst{latency_bottleneck}, \kst{latency_unknown_arch} \\
\bottomrule
\end{tabularx}
\caption{The analyzer's ten analyses. The occupancy model is the min-wins
formula from the occupancy library; the latency estimate is explicitly
approximate.}
\label{tab:c1checks}
\end{table}

The occupancy estimate is:
\[
  \text{blocks} = \min\left(
  \left\lfloor \frac{\text{threads}_{\max}}{\text{threads per block}}\right\rfloor,
  \left\lfloor \frac{\text{warps}_{\max}}{\text{warps per block}}\right\rfloor,
  \left\lfloor \frac{\text{smem}_{\max}}{\text{smem per block}}\right\rfloor,
  \text{blocks}_{\text{regs}},
  \text{blocks}_{\text{hard}}
  \right),
\]
and the binding term is reported as the limiter. The bank-conflict metric
uses $32$ banks of $4$ bytes: the bank index is
$(\text{word} \div 4) \bmod 32$, and the conflict way is lanes divided by
distinct banks. Coalescing uses the 128-byte-sector model and reports a
byte-efficiency ratio; it does not claim a hardware counter.

\subsection{Suggestion rules}
\label{sec:c1-suggest}

The suggestion engine owns \kst{report.suggestions}. It clears the list and
emits at most one suggestion per triggering finding, in a fixed priority
order, with an evidence text and a \kst{triggered_by} link.

\begin{table}[htbp]
\centering
\small
\begin{tabularx}{\textwidth}{@{}l X l@{}}
\toprule
Rule & Trigger and recommendation & Evidence \\
\midrule
\kst{tile_k_32} & Bank conflict of way $\ge 2$; set the shared tile to 32 words. & Current stride, banks, shared bytes \\
\kst{split_accumulator} & Register limiter and at least 16 FP32 FMAs; split the accumulator. & Registers, FMA count, fragment size \\
\kst{hoist_predicate} & Divergent branch inside an inner loop; hoist the predicate. & Branch site, predicate dependence \\
\kst{vectorize_coalescing} & Coalescing efficiency below 1.0 and no bank conflict; target 128\,B per warp. & Efficiency, sectors, vector width \\
\bottomrule
\end{tabularx}
\caption{Suggestion rules. The priority order is the table order. The engine
drops instruction-mix informational findings so a suggestion is tied to an
actionable warning.}
\label{tab:c1suggest}
\end{table}

\subsection{Report and golden result}
\label{sec:c1-evidence}

The report is a JSON object with \kst{kernel}, \kst{metrics}, \kst{limiter},
\kst{warnings}, and \kst{suggestions}; the schema is summarized in
Appendix~\ref{app:report-schema}. The committed golden fixture is a
fused-attention-style CUDA kernel whose PTX targets sm\_86:
occupancy $9 \text{ blocks} \times 2 \text{ warps} / 48 = 0.375$, limiter
\kst{registers}, 102 registers per thread, coalescing efficiency $0.5$,
a $68$-word shared stride that maps to banks $\{0,4,\dots,28\}$ with way $4$,
two \kst{bar.sync} instructions, 1024 FP32 FMAs, zero spills, and three
warnings plus three suggestions. The sm\_86 fixture is a parser corpus
artifact; it is independent of the local \kst{nvcc} capability pin.

% ===========================================================================
\section{C2: correctness-gated autotuner}
\label{sec:c2}

\kst{kernscope-tune} searches a kernel family's configuration grid for the
fastest variant that passes a bf16 golden gate. It reports the winner with
roofline and occupancy evidence, the full variant list, and a reproducibility
bundle. It is not a blind search: infeasible configurations are rejected
before compilation, and wrong kernels are never benchmarked.

\subsection{Families and search spaces}
\label{sec:c2-space}

\begin{table}[htbp]
\centering
\small
\begin{tabularx}{\textwidth}{@{}l X@{}}
\toprule
Family & Auto grid \\
\midrule
\kst{matmul} & $(BM,BN) \in \{(64,64),(128,128)\} \times BK \in \{16,32\} \times
warps \in \{1,2,4\} \times stages \in \{1,2\}$: 24 configs. The full grid
expands to 4608 combinations over $BK$ 8--64, stages 1--4, vector width 1--4,
XOR swizzle, and unroll 1--4. \\
\kst{bias_gelu}, \kst{softmax}, \kst{layernorm}, \kst{reduction} &
block $\in \{128,256,512,1024\}$ $\times$ vector width $\in \{1,2,4,8\}$;
softmax and layernorm add a one-pass/two-pass mode, reduction adds
strided/tree. \\
\kst{flash_attention} & $Br \times Bc \in \{32,64\}^2$ $\times$ stages
$\{1,2\}$ $\times$ vector width $\{1,2,4\}$ $\times$ shared-memory layout
(padding/XOR) $\times$ causal on/off. \\
\bottomrule
\end{tabularx}
\caption{Tuner families and their search spaces.}
\label{tab:c2space}
\end{table}

Each family owns a template generator that validates membership in the search
space and checks the \fpath{libs/arch} budgets (registers, shared memory,
warp/block divisibility, vector tiling) before compilation. A rejection is a
reason string, not a compiler error.

\subsection{The pipeline}
\label{sec:c2-pipeline}

\begin{figure}[htbp]
\centering
\resizebox{\textwidth}{!}{% Tuner pipeline: template path and capture path.
\begin{tikzpicture}[
    font=\small,
    box/.style={draw=ink!35, rounded corners=3pt, align=center, text width=2.3cm,
                minimum height=1.3cm, inner sep=3pt, font=\scriptsize, fill=paper},
    flow/.style={-{Stealth[length=2mm]}, line width=0.8pt, draw=ink!60},
    dflow/.style={-{Stealth[length=2mm]}, line width=0.7pt, dashed, draw=ink!55},
    lbl/.style={font=\scriptsize\itshape, text=ink!60},
  ]

  % Template path.
  \node[box, fill=kgreenfill, draw=kgreen] (family) at (0,2.1)
    {\textbf{family + shape}\\matmul, bias\_gelu,\\softmax, layernorm,\\reduction, attention};
  \node[box, fill=paper] (grid) at (3.1,2.1)
    {\textbf{config grid}\\tiles, warps,\\stages, vector width};
  \node[box, fill=kgoldfill, draw=kgold] (screen) at (6.2,2.1)
    {\textbf{pre-screen}\\register, smem,\\occupancy budgets};
  \node[box, fill=kbluefill, draw=kblue] (compile) at (9.3,2.1)
    {\textbf{compile}\\hipcc gfx1101\\or nvcc sm\_80/90};

  \draw[flow] (family) -- (grid);
  \draw[flow] (grid) -- (screen);
  \draw[flow] (screen) -- (compile);

  % The rest wraps to a second row.
  \node[box, fill=kgreenfill, draw=kgreen] (golden) at (0,-1.1)
    {\textbf{golden gate}\\bfloat16 cosine,\\max-rel, PCC};
  \node[box, fill=korangefill, draw=korange] (bench) at (3.1,-1.1)
    {\textbf{benchmark}\\warmup, L2 flush,\\median + p20/p80};
  \node[box, fill=kvioletfill, draw=kviolet] (select) at (6.2,-1.1)
    {\textbf{select}\\median, then p20,\\then p80};
  \node[box, fill=paper] (cache) at (9.3,-1.1)
    {\textbf{cache + db}\\arch:family:shape\\JSONL records};

  \draw[flow] (compile.south) -- ++(0,-1.0) -| (golden.north);
  \draw[flow] (golden) -- (bench);
  \draw[flow] (bench) -- (select);
  \draw[flow] (select) -- (cache);

  % Capture path.
  \node[box, fill=kgoldfill, draw=kgold] (bundle) at (0,-4.3)
    {\textbf{capture bundle}\\binary + buffers\\+ launch};
  \node[box, fill=paper] (launch) at (3.1,-4.3)
    {\textbf{launch configs}\\captured grid,\\block search};
  \node[box, fill=kgreenfill, draw=kgreen] (replay) at (6.2,-4.3)
    {\textbf{bitwise replay}\\output hash\\must match};
  \node[box, fill=korangefill, draw=korange] (cbench) at (9.3,-4.3)
    {\textbf{benchmark}\\same harness,\\same quantiles};

  \draw[flow] (bundle) -- (launch);
  \draw[flow] (launch) -- (replay);
  \draw[flow] (replay) -- (cbench);
  \node[lbl, above] at (3.1,-3.65) {\texttt{--from-capture}};

  % Shared selection.
  \draw[dflow] (cbench.south) -- ++(0,-0.75) -| (select.south);

\end{tikzpicture}
}
\caption{The tuner pipeline. The template path searches source variants;
the capture path searches launch configurations over a captured binary. Both
use the same benchmark and selection rules.}
\label{fig:tuner}
\end{figure}

A candidate moves through five stages:
\begin{enumerate}
  \item \textbf{Cache hit.} A config already in the tunecache is replayed
        without compiling.
  \item \textbf{GPU-free pre-screen.} The generator validates the config, then
        \kst{libs/occupancy} rejects zero-resident-block configs and configs
        below \kst{--min-occupancy}.
  \item \textbf{Compile.} The family generator emits a single-source
        \kst{.hip} file; \kst{hipcc --genco} compiles it for gfxXXXX, or
        \kst{nvcc -c} for sm\_XX.
  \item \textbf{Golden gate.} The kernel runs on the live HIP device through
        the template runner and the bf16 reference. The gate passes only when
        cosine $\ge 0.999$, maximum relative error $\le 10^{-2}$, and PCC
        $\ge 0.99$. A failure discards the variant; it is never timed.
  \item \textbf{Benchmark and select.} \fpath{libs/bench} warms up, flushes
        L2, adapts the trial count, and reports median and p20/p80. The best
        config wins by median, with p20, p80, and enumeration order as
        tie-breakers.
\end{enumerate}

The gate shape is the family's golden shape (64$^3$ for matmul), deliberately
decoupled from the benchmark shape. The cache is keyed
\kst{gpu_arch:family:shape_hash} and stamped with the minimum-occupancy floor
that produced it, so a floor-1.0 entry never replays for a floor-0.0 run.
\kst{--no-cache} skips the read but still writes the new entries.

\subsection{Outputs and reproducibility}
\label{sec:c2-output}

\kst{--json} emits the best config and timings, the baseline (the first
enumerated measured config), the speedup over that baseline, every measured
variant sorted by median, the pipeline counts, the roofline position and bound,
and the winning occupancy and limiter. \kst{--text} is the owner-facing
one-screen form. \kst{--html} is a self-contained variant comparison with the
model embedded under \kst{report-data} and the evidence under
\kst{tune-evidence}.

\kst{--emit-repro} writes a standalone directory with \kst{run.sh},
\kst{README.md}, and the winning config. The script regenerates the source,
repeats the exact compile command, and re-runs the golden gate at the winning
geometry. The nixpkgs revision from \kst{devenv.lock} is embedded and checked
against \kst{KERNSCOPE_EXPECT_NIXPKGS_REV} as a warning. The repro path needs
a live HIP device and exists for matmul plus the four elementwise families;
\kst{flash_attention} has no generator CLI and is rejected with a clear
message.

\subsection{Capture-tune mode}
\label{sec:c2-capture}

A capture bundle contains a binary, arguments, and buffers, but no source.
Template tuning cannot consume it; the launch configuration is the only
searchable degree of freedom. \kst{kernscope-tune --from-capture} restores the
captured launch, enumerates block sizes at the captured grid, pre-screens the
limits (threads per block $\le 1024$, grid x $\le 2^{31}-1$, grid y/z
$\le 65535$, shared memory within the architecture budget), replays each
survivor bitwise, and times only the survivors. The captured baseline is
candidate 1, so \kst{best.speedup} is always at least 1.0. The record key is
\kst{(arch, sha256(kernel name), sha256(bundle identity))}; the binary and
buffers stay byte-identical to the capture.

\subsection{Measured evidence}
\label{sec:c2-evidence}

In the demo, capture-tune explored four candidates over the captured
layernorm bundle. All four validated bitwise; the best block was 256 with a
median of $0.01206$\,ms against a $0.01208$\,ms baseline, a speedup of
$1.0017\times$. The verifier-gated template run on bias\_gelu at
$4096 \times 4096$ enumerated 13 variants, compiled 12, and gated one before
compilation. The best config (block 256, vector width 4) ran at
$0.17184$\,ms against a $0.20332$\,ms baseline, a $1.1832\times$ speedup,
occupancy 1 (limiter registers), arithmetic intensity 2.333, and roofline
position 0.0845 in the memory-bound class.

The dashboard records a separate measured matmul search with 15 variants, a
best median of $415.324\,\mu$s, and a worst median of $3.106$\,ms. The
per-run JSONL for that search is not committed, so the report quotes the
dashboard value and labels it as such rather than reconstructing a raw record.

% ===========================================================================
\section{C3: static verifier}
\label{sec:c3}

\kst{kernscope-verify} compiles CUDA or HIP source to LLVM IR at \kst{-O1}
and applies eight may-analyses. It does not prove a kernel correct; it proves
specific bounds and ordering facts when it can, and it refuses approval when
it cannot.

\subsection{Ingestion and launch tiers}
\label{sec:c3-ingest}

By default the host and device halves are compiled separately and linked into
one module, so the launch site is visible. The host pass supplies the launch
configuration through the \kst{__cudaPushCallConfiguration} or
\kst{cudaLaunchKernel} call. HIP uses clang 21 with \kst{--hip-path}.
\kst{--device-only} is a degraded mode: launch dimensions are lost unless the
kernel carries \kst{__launch_bounds__}.

Launch dimensions resolve into three tiers:
\begin{enumerate}
  \item \textbf{Exact}: constant dimensions at the launch site, giving exact
        ranges for \kst{threadIdx} and \kst{blockIdx}.
  \item \textbf{UpperBound}: \kst{__launch_bounds__} or \kst{nvvm.maxntid}
        provides a sound upper bound on \kst{blockDim}.
  \item \textbf{Symbolic}: runtime dimensions produce a \kst{SYMBOLIC_DIMS}
        verdict and a runtime-check suggestion, never a claim of safety.
\end{enumerate}

\begin{figure}[htbp]
\centering
\resizebox{\textwidth}{!}{% Verifier pipeline: source to verdicts to approved IR.
\begin{tikzpicture}[
    font=\small,
    box/.style={draw=ink!35, rounded corners=3pt, align=center, text width=2.45cm,
                minimum height=1.35cm, inner sep=3pt, font=\scriptsize, fill=paper},
    checks/.style={draw=kblue, rounded corners=3pt, align=center, inner sep=4pt,
                   font=\scriptsize, fill=kbluefill},
    flow/.style={-{Stealth[length=2mm]}, line width=0.8pt, draw=ink!60},
    dflow/.style={-{Stealth[length=2mm]}, line width=0.7pt, dashed, draw=ink!55},
    lbl/.style={font=\scriptsize\itshape, text=ink!60},
  ]

  % Ingestion and launch-config tiers.
  \node[box, fill=kgreenfill, draw=kgreen] (src) at (0,1.8)
    {\textbf{CUDA / HIP}\\source};
  \node[box, fill=paper] (clang) at (3.25,1.8)
    {\textbf{clang}\\host + device\\into one module};
  \node[box, fill=paper] (ir) at (6.5,1.8)
    {\textbf{LLVM IR}\\at \texttt{-O1}\\SCEV + LVI ranges};
  \node[box, fill=kgoldfill, draw=kgold] (tiers) at (9.75,1.8)
    {\textbf{launch tiers}\\Exact / UpperBound\\/ Symbolic};

  \draw[flow] (src) -- (clang);
  \draw[flow] (clang) -- (ir);
  \draw[flow] (ir) -- (tiers);

  % Check registry and verdicts.
  \node[checks] (checks) at (4.5,-1.9)
    {\textbf{check registry (8 classes)}\\[3pt]
     \begin{tabular}{@{}ll@{}}
       1 global OOB & 2 shared bounds\\
       3 barrier & 4 race\\
       5 pointer & 6 divergent barrier\\
       7 uninitialised & 8 address space
     \end{tabular}\\[3pt]
     {\scriptsize per-function budget 65536}};
  \node[box, fill=kvioletfill, draw=kviolet] (verdicts) at (8.4,-1.9)
    {\textbf{verdict document}\\per class: \texttt{SAFE}\\
     or \texttt{MAY\_*,\,UNKNOWN\_*}\\
     + source locations};
  \node[box, fill=kgreenfill, draw=kgreen] (approved) at (11.6,-1.9)
    {\textbf{approved IR}\\only if every\\verdict is \texttt{SAFE}};
  \node[box, fill=kredfill, draw=kred] (reject) at (8.4,-4.1)
    {\textbf{reject}\\exit 1; no\\\texttt{.verified.ll}};

  \draw[flow] (ir.south) -- (checks.north);
  \draw[flow] (tiers.south) -- (checks.north east);
  \draw[flow] (checks.east) -- (verdicts.west);
  \draw[flow] (verdicts.east) -- (approved.west);
  \draw[flow, kred] (verdicts.south) -- (reject.north);

  % The eBPF design analogy.
  \node[font=\scriptsize\itshape, text=ink!55, anchor=west, text width=12.6cm] at (0,-5.25)
    {The UX follows the Linux eBPF verifier: reject with a line-level reason,
     track types and ranges, bound the explored states, and degrade to a
     non-proof verdict instead of hanging or claiming safety.};
\end{tikzpicture}
}
\caption{The verifier pipeline. The launch tiers feed the diagnostic engine;
the check registry produces per-class verdicts; approved IR is written only
when no error-class verdict remains.}
\label{fig:verify}
\end{figure}

\subsection{The eight check classes}
\label{sec:c3-checks}

\begin{table}[htbp]
\centering
\small
\begin{tabularx}{\textwidth}{@{}c l l X@{}}
\toprule
ID & Class & Plugin & Verdict vocabulary \\
\midrule
1 & Global out-of-bounds & \kst{checks/oob} & \kst{SAFE}, \kst{MAY_OOB}, \kst{SYMBOLIC_DIMS} \\
2 & Shared-memory bounds & \kst{checks/oob} & \kst{SAFE}, \kst{MAY_OOB}, \kst{SYMBOLIC_DIMS} \\
3 & Barrier correctness & \kst{checks/barrier} & \kst{SAFE}, \kst{BARRIER_MISMATCH} \\
4 & Race-prone shared writes & \kst{checks/race} & \kst{SAFE}, \kst{MAY_RACE}, \kst{UNKNOWN_RACE} \\
5 & Illegal pointer arithmetic & \kst{checks/ptr} & \kst{SAFE}, \kst{PTR_OOB} \\
6 & Divergent barriers & \kst{checks/barrier} & \kst{SAFE}, \kst{DIVERGENT_BARRIER} \\
7 & Uninitialized reads & \kst{checks/ptr} & \kst{SAFE}, \kst{MAY_UNINIT_READ} \\
8 & Invalid address-space casts & \kst{checks/ptr} & \kst{SAFE}, \kst{INVALID_ADDRSPACE_CAST} \\
\bottomrule
\end{tabularx}
\caption{The eight check classes and their stable IDs. Engine-degraded verdicts
use IDs at or above 100: symbolic dims (100), SCEV reconstruction failed (101),
and budget exceeded (102).}
\label{tab:c3checks}
\end{table}

The design follows the Linux eBPF verifier's user experience
\citep{ebpfVerifier}: reject with a concrete line-level reason, track types and
ranges, cap the explored states, and degrade to a non-proof verdict rather
than hanging. Every non-\kst{SAFE} verdict carries diagnostics with a source
location, the enclosing kernel, the offending instruction, and an evidence
path where one exists. The per-function explored-state budget defaults to
65536 and is settable with \kst{--budget}.

\subsection{Verdict semantics and approved IR}
\label{sec:c3-verdict}

\kst{SAFE} is the only approval-class verdict. Every other verdict, including
the engine-degraded ids, means no bound was proven. The CLI exits 1 and writes
no approved IR. If the error-class count is zero, the verifier projects the
verified module onto its device code: it keeps the kernels and reachable
device helpers, drops host functions, and rewrites the global constructor and
used lists. The result is \kst{<base>.verified.ll}. The projection matters for
the handoff because the proofs apply to device code and the AMDGPU backend
rejects host \kst{addrspace(0)} allocas.

\subsection{Soundness scope}
\label{sec:c3-soundness}

All eight checks are may-analyses at \kst{-O1}. In scope: global and shared
accesses whose indices are bounded by SCEV and the launch-dimension lattice;
barrier counts and divergent conditional barriers; intra-kernel shared-memory
access pairs partitioned by barrier phases; pointer arithmetic, uninitialized
reads, and address-space casts. Out of scope: data-dependent indices (reported
\kst{MAY_OOB} or \kst{UNKNOWN_RACE}, never \kst{SAFE}), runtime-parameter
allocations (an index-bounded \kst{SAFE} carries an explicit note), atomics,
and host-side allocation tracking.

The known false positives include identical thread-indexed expressions that a
value-numbering pass could merge, disjoint arithmetic that SCEV cannot prove
disjoint, unbounded thread indices that overlap everything, loop-carried
phases, and an unknown-range thread access paired with a uniform address. The
known false negatives include data-dependent indices (named \kst{UNKNOWN_RACE},
not proven safe), uniform-address value races, lone data-dependent stores,
divergent barriers still treated as separators, and \kst{__syncwarp},
\kst{__threadfence_block}, and atomics not modeled as synchronization.
Inter-block races are out of scope. The future path is HGRD-style
dominance-based barrier intervals, Faial-style behavioral types, and a
GPUVerify-style two-thread reduction; none is implemented here.

\subsection{Corpus and evidence}
\label{sec:c3-evidence}

The corpus has 19 source files and 34 consolidated manifest entries across
the four check plugins. The lit gate runs the CLI \kst{--json} over every
corpus kernel and diffs the verdicts against the manifest read from disk; the
CLI never reads the manifest, so the test is not self-reporting. The
dashboard summarizes the projection as 18 \kst{SAFE}, 4 \kst{MAY_OOB},
4 \kst{PTR_OOB}, 2 \kst{MAY_RACE}, 2 \kst{MAY_UNINIT_READ},
1 \kst{DIVERGENT_BARRIER}, 1 \kst{BARRIER_MISMATCH}, 1 \kst{UNKNOWN_RACE},
and 1 \kst{INVALID_ADDRSPACE_CAST}. In the demo, the verifier gate removed one
risky bias\_gelu variant before compilation while 12 others compiled and
benchmarked.

% ===========================================================================
\section{C4: capture and replay}
\label{sec:c4}

\kst{kernscope-capture} records what a kernel launch observably did and lets
a later process replay it. The recorded bundle is the input to the analyzer,
the capture-tune search, and a standalone reproducer. The central constraint is
VA fidelity: the replay must reconstruct the original pointer graph, not just
the bytes.

\subsection{Mechanisms}
\label{sec:c4-mechanisms}

\Cref{tab:capture-mechanisms} lists the three mechanisms. All three share the
\fpath{libs/capture-core} session model and bundle writer, so the reader does
not need to know which one produced a bundle. Vendor selection is
\kst{--vendor}, then \kst{$KERNSCOPE_VENDOR}, then the first capture shared
object next to the CLI (HIP preferred), then an error.

\subsection{Bundle format v1}
\label{sec:c4-format}

\begin{figure}[htbp]
\centering
\resizebox{\textwidth}{!}{% Capture bundle v1: one folder, one event log, one VA snapshot.
\begin{tikzpicture}[
    font=\small,
    stage/.style={draw, rounded corners=3pt, align=center, text width=2.3cm,
                  minimum height=1.1cm, inner sep=3pt, font=\scriptsize, fill=paper},
    part/.style={draw=ink!30, rounded corners=2pt, align=center, text width=1.75cm,
                 minimum height=1.0cm, inner sep=3pt, font=\scriptsize, fill=ink!3},
    flow/.style={-{Stealth[length=2mm]}, line width=0.8pt, draw=ink!60},
    note/.style={font=\scriptsize\itshape, text=ink!60, align=left},
  ]

  % The bundle folder.
  \draw[draw=kviolet, fill=kvioletfill!45, rounded corners=4pt, line width=0.8pt]
    (0.2,-1.9) -- (8.8,-1.9) -- (8.8,2.1) -- (4.1,2.1) -- (3.5,2.7) -- (0.2,2.7) -- cycle;
  \node[font=\small\bfseries, text=kviolet, anchor=west] at (0.35,2.3) {bundle.v1/};

  \node[part, fill=paper] (man) at (1.75,0.75)
    {\texttt{manifest.json}\\{\tiny format, vendor, arch,\\kernel, hashes}};
  \node[part, fill=paper] (events) at (1.75,-0.95)
    {\texttt{events.json}\\{\tiny launches, args,\\memcpys, syncs}};
  \node[part, fill=paper] (blobs) at (4.55,0.75)
    {\texttt{blobs.json}\\{\tiny VA ranges, sizes,\\provenance}};
  \node[part, fill=paper] (snap) at (4.55,-0.95)
    {\texttt{blobs/*.hsaco}\\{\tiny code objects and\\buffer snapshots}};
  \node[part, fill=paper] (repro) at (7.25,0.0)
    {\texttt{reproducer/}\\{\tiny standalone C++\\project}};

  % Record and replay around the folder.
  \node[stage, fill=korangefill, draw=korange] (record) at (-2.0,0.75)
    {\textbf{record}\\running binary};
  \node[stage, fill=kgreenfill, draw=kgreen] (replay) at (11.3,0.75)
    {\textbf{replay}\\restore + launch};
  \node[stage, fill=kbluefill, draw=kblue] (check) at (11.3,-1.55)
    {\textbf{validate}\\bitwise hash};
  \draw[flow] (record) -- (0.2,0.75);
  \draw[flow] (8.8,0.75) -- (replay);
  \draw[flow] (replay) -- (check);
  % The patch table is the subtle part of VA fidelity.
  \node[note, anchor=north west, text width=4.8cm] at (0.35,-2.12)
    {embedded pointers are recorded in a patch table and rewritten to fresh
     virtual addresses on replay; the output hash must match the captured
     reference exactly};
\end{tikzpicture}
}
\caption{Bundle v1. The manifest names the vendor and mechanism; the event log
records API calls, allocations, launches, copies, memsets, streams, modules,
and parameter metadata; the blob index carries a SHA-256 and a patch table for
every snapshot.}
\label{fig:capture}
\end{figure}

A v1 bundle is a directory:
\begin{lstlisting}
<name>.kcap/
  manifest.json
  events.json
  blobs.json
  blobs/b0001.bin
\end{lstlisting}

\kst{manifest.json} must have \kst{format_id="kernscope-capture"},
\kst{schema_version="1.0"}, \kst{vendor} in \{nvidia, amd, mock\},
\kst{capture_mechanism} in \{cupti, preload, hip-sdk\},
\kst{snapshot_policy="va-faithful"}, and a \kst{size_governance} object with a
default maximum snapshot of 1\,GiB. A total over the cap stops further blob
writes, sets \kst{aborted=true}, records a reason, and leaves a loadable
partial bundle.

\kst{events.json} is an ordered log. The record types are API counts,
allocation and free, unknown free, launch, copy, memset, stream, fatbin,
cubin, and parameter info. A launch carries the kernel name, the launch
mechanism, grid/block/shared memory, stream and graph handles, and a
per-parameter argument table. Each argument is a pointer or scalar, with a
size, a byte-exact host copy, and the original virtual address for pointer
slots.

\kst{blobs.json} indexes the snapshots. Each blob has an id, file, role, VA,
size, SHA-256, and patch table. Roles are kernel binary, input, output,
argument, and state. The reader recomputes every hash on read and rejects a
mismatch. The patch table is the mechanism that makes replay portable: every
8-byte little-endian value in a blob that equals a known allocation VA is
recorded as a patch, and replay rewrites only the entries present in the
allocation remap. A \kst{__device__} global that happens to contain another
global's address is out of scope.

\subsection{Replay, validation, and benchmark}
\label{sec:c4-replay}

The replay model resolves buffers from allocations, H2D inputs, D2H markers,
and v1 snapshots, then restores the binary, the argument plan, and the launch
geometry. HIP/HSA assigns virtual addresses, and the current runtime wrapper
does not expose CUDA VMM placement, so replay remaps on both vendors. It
allocates fresh buffers, computes old-to-new VA remapping, applies the patch
tables before upload, rewrites pointer argument slots through the same table,
and launches the captured grid, block, and shared memory.

Validation is bitwise by default: byte compare plus SHA-256, with a hash diff
on failure. Tolerance mode interprets outputs as bf16 chunks upcast to fp32
and passes only when cosine $\ge 0.999$, maximum relative error
$\le 10^{-2}$, and PCC $\ge 0.99$. References come from v1 output snapshots
or from \kst{--ref} in output-buffer order; a missing reference is an error,
not a pass.

\kst{bench} performs one validation launch before timing and refuses to time
an invalid replay unless \kst{--no-validate} is given. It uses ten repetitions
by default, an L2 flush, event timing, and median plus p20/p80. Each run
writes one \kst{libs/db} record keyed by architecture and kernel hash, with
\kst{provenance.mode="measured"}.

\subsection{Reproducer and size governance}
\label{sec:c4-repro}

\kst{emit-repro} writes a standalone plain-CMake C++ project with the kernel
binary, inputs, a \kst{main.hip} or \kst{main.cu}, a \kst{CMakeLists.txt}
pinned to the manifest vendor and architecture, and a README. It re-points
pointer arguments at fresh buffers and launches the captured geometry.
Size-unknown buffers become deterministic synthetic inputs, marked so a reader
does not mistake them for bit-exact data. The NVIDIA reproducer is
compile-only.

\subsection{CUDA/HIP field mapping}
\label{sec:c4-mapping}

The mapping contract compares what the two vendors can both express, after
normalization. Kernel names are compared after stripping the \kst{.kd} suffix
and demangling. Grid, block, and shared memory compare directly. The default
stream normalizes to 0; non-default streams get a 1-based id. Pointers compare
as 1-based allocation ordinals, not raw VAs. Scalars compare by size and
byte-identical value; input blobs compare by SHA-256 after H2D call order.
Allocations compare as a sorted multiset of sizes; copies compare by kind
(H2D, D2H, D2D). Binary blobs are compared per vendor only; a cubin is never
compared with an hsaco. API counts, raw pointers before normalization, and
timestamps are not compared.

\subsection{Measured evidence}
\label{sec:c4-evidence}

The demo captured the hot layernorm kernel
\kst{_Z15layernorm_fused}\allowbreak\kst{PKtS0_S0_Ptfi.kd} into a v0
text bundle. The v1 layout is the current reader contract and the
capture-tune path accepts it, but the record CLI in this tree still writes
the v0 form; the gap is flagged in \Cref{sec:limitations-list}. The launch
was at grid
$1024 \times 1 \times 1$ and block $128 \times 1 \times 1$. The bundle has
three blobs. Replay reproduced the output bitwise; the reference SHA-256 is
\kst{a0027ed9a8ef2e3453aeb194aa64108c}\allowbreak\kst{1f300a87fbfa292cca5222a725449eec}.
The capture-tune stage then explored four block sizes over the same bundle,
validated all four bitwise, and selected block 256. This is the strongest
cross-tool evidence in the report: the analyzer's stage-3 findings and the
tuner's stage-4 candidates both refer to the same captured code object and
launch.

% ===========================================================================
\section{C5: LLVM passes}
\label{sec:c5}

\kst{kernscope-pass} is an out-of-tree LLVM 21 pass plugin. It is loaded by
both driver paths, \kst{opt -load-pass-plugin} and
\kst{clang -fpass-plugin}, and it uses the New Pass Manager's
\kst{PassPluginLibraryInfo} entry point. The \kst{kernscope-pass-report} CLI
runs a fixed showcase pipeline and measures the before/after delta with the
analyzer.

\subsection{Plugin structure}
\label{sec:c5-plugin}

A pass module subclasses \kst{llvm::}\allowbreak\kst{PassInfoMixin} and
registers a callback with \kst{REGISTER_PASS_}\allowbreak\kst{CALLBACK} at
file scope. The callback receives the
\kst{PassBuilder} of the loading process and wires the module into the
pipeline. The registry is a process-global singleton; the plugin entry point
does not need to know the module names in advance. The lit suite exercises
both \kst{opt} and \kst{clang} loading so an ABI change fails at the loader
rather than at the first transformed kernel.

\subsection{The passes}
\label{sec:c5-passes}

\begin{table}[htbp]
\centering
\small
\begin{tabularx}{\textwidth}{@{}l@{\hspace{1em}}X@{}}
\toprule
Pass & Transformation or analysis \\
\midrule
\kst{kernscope-promote} & Global-to-shared promotion for tiles reused in a loop. \\
\kst{kernscope-tile} & Loop tiling over constant-trip, unit-stride \kst{\{s,+,1\}} windows. \\
\kst{kernscope-vectorize} & Rewrites scalar load/store chains into vector operations when alignment and contiguity permit. \\
\kst{kernscope-pad} & Shared-memory padding to break a known bank-conflict stride. \\
\kst{kernscope-rle} & Redundant-load elimination and address simplification. \\
\kst{kernscope-coalesce} & Analysis of per-access coalescing class; no rewrite. \\
\bottomrule
\end{tabularx}
\caption{The pass modules. The coalescing analysis is shared with the
analyzer's 128-byte-sector model so before/after numbers use the same
definition.}
\label{tab:c5passes}
\end{table}

The default showcase pipeline is
\kst{kernscope-tile},\allowbreak
\kst{kernscope-promote},\allowbreak
\kst{kernscope-vectorize},\allowbreak
\kst{kernscope-pad},\allowbreak
\kst{kernscope-coalesce}.
Not every pass fires on every kernel; a no-transform result is a valid
outcome and is recorded as such. The tile pass, for example, declines a
layernorm loop whose induction variable is not a canonical \kst{\{0,+,1\}}
phi.

\subsection{The delta report}
\label{sec:c5-delta}

\kst{kernscope-pass-report --report-pipeline <gemm|layernorm|softmax>} runs
clang to before-IR, the plugin pipeline to after-IR, and \kst{llc} to
before/after PTX. It then invokes \kst{kernscope-analyze} on both PTX files
and emits a delta report with register estimate, shared-memory bytes,
coalescing class, bank-conflict way, vector width, and vectorized global
loads. Every metric is labeled improved, unchanged, or not applicable. The
report's provenance is \kst{projected}: the IR and PTX are compile-only
artifacts, and the note says the measured gfx1101 analog lives in the
\kst{libs/db} records.

On layernorm, vector width goes from 1 to 4, vectorized global loads from 0 to
12, uncoalesced accesses from 64 to 16, and the register estimate stays at
175. The pass report records 12 vector chains and 48 vectorized loads. gemm
and softmax are correctly left no-transform, and all three kernels compile to
PTX with exit 0 after the pipeline. The tile pass's refusal is recorded with
its reason rather than hidden.

% ===========================================================================
\section{C6: profiler}
\label{sec:c6}

\kst{kernscope-profile} answers three questions from a running binary: which
kernels ran, how long each took, and what the compiler-visible facts suggest
about the bottleneck. It does not require hardware counters. It traces, it
enriches, it ranks, and it diagnoses from the activity records and the
architecture database.

\subsection{Pipeline}
\label{sec:c6-pipeline}

\begin{enumerate}
  \item \textbf{Trace.} The target runs under \kst{LD_PRELOAD} of the
        rocprofiler-sdk tracer; each kernel dispatch and HIP API call becomes a
        JSONL record. On NVIDIA the equivalent path is the CUPTI activity
        tracer, which is compile- and unit-tested locally and rental-gated for
        live use.
  \item \textbf{Enrich.} Each kernel record receives occupancy and limiter
        estimates from \fpath{libs/occupancy}, code-object VGPR/SGPR/LDS
        facts where available, and explicit degradation markers where not.
  \item \textbf{Rank.} Kernels are grouped by name and ranked by total time
        (then median, then name); the top ten are kept.
  \item \textbf{Diagnose.} The rule engine classifies memory, latency, and
        compute behavior for known families, reports the occupancy limiter,
        and degrades to \kst{insufficient_evidence} rather than inventing a
        cause.
  \item \textbf{Report.} \kst{--json} emits the model, \kst{--text} the
        owner-facing table, and \kst{--html} a self-contained page with the
        embedded JSON.
\end{enumerate}

\subsection{Diagnostics}
\label{sec:c6-diagnostics}

A \kst{memory_bound} hint is emitted when the workload's arithmetic intensity
is below the measured ridge point. A \kst{latency_bound} or
\kst{compute_bound} hint follows the same roofline position and the
instruction mix. An occupancy limiter hint is emitted when occupancy is below
1.0; the limiter is the binding resource. Sub-hints distinguish register
bloat, shared-memory pressure, and a block too small to fill the machine.
Every hint carries its evidence: arithmetic intensity, bytes moved, ridge
point, occupancy, registers, or the analyzer finding id when the analyzer was
merged into the pipeline.

\subsection{Measured evidence}
\label{sec:c6-evidence}

Tracing the layernorm app produced three kernels and 19.6\,$\mu$s of wall
time. The hot kernel ran for 9.2\,$\mu$s with 24 registers per thread,
occupancy 1 (limiter registers), 1024 bytes of shared memory, grid
$1024 \times 1 \times 1$, and block $128 \times 1 \times 1$. The profiler
emitted a memory-bound hint with arithmetic intensity 1.375, a bytes-moved
estimate of 1048576, and a ridge point of 27.614. The hint's
\kst{analyzer_cross_refs} in the aggregated report point at the stage-3
analyzer findings, including \kst{possibly_divergent_branch}.

% ===========================================================================
\section{Bench, dashboard, and the evidence layer}
\label{sec:dashboard}

\subsection{Benchmark methodology}
\label{sec:bench-method}

\kst{libs/bench} is the one timing method in the project. It warms up the
kernel, flushes L2 between trials, and adapts the trial count to a time budget
while keeping a minimum sample count. It reports the median and the p20/p80
quantiles, not a mean, so a single slow launch does not move the result. The
same harness times the peak microbenchmarks, the tuner variants, and the
capture-replay benchmark. A correctness gate runs before timing in the
tuner and in capture replay; the peak microbenchmarks are their own reference.

\subsection{Measured machine peaks}
\label{sec:peaks}

The peaks in \Cref{tab:peaks} drive the roofline engine. The fp32 figure is
the register-resident FFMA loop; the DRAM read figure uses a 512\,MiB array;
the LDS figure is a read-plus-write loop. The device-to-device figure is
reported as a range over the accepted runs. Clocks were recorded with
\kst{rocm-smi} but not locked, and the machine was not isolated from
development activity. The numbers are therefore a measured description of
this machine on 2026-08-14, not a datasheet claim.

\subsection{Dashboard sections}
\label{sec:dashboard-sections}

The dashboard aggregates six sections. Each section has a title, an
availability flag, a status, a provenance object, a summary, and a list of
evidence paths.

\begin{table}[htbp]
\centering
\small
\begin{tabularx}{\textwidth}{@{}l l X@{}}
\toprule
Section & Mode & Summary at consolidation \\
\midrule
\kst{analyzer} & projected & Golden occupancy 0.375, 102 registers, limiter registers, 3 warnings, 3 suggestions. \\
\kst{tune} & measured & 15 matmul variants; best 415.324\,$\mu$s, worst 3.106\,ms. \\
\kst{verify} & projected & 34 entries, 19 files, 8 classes; 18 SAFE. \\
\kst{capture} & measured & Grid $1024\times1\times1$, block $128\times1\times1$, 3 blobs, bitwise PASS. \\
\kst{profile} & measured & 3 kernels, 19.6\,$\mu$s wall, hot 9.2\,$\mu$s, memory-bound hint. \\
\kst{pass} & projected & Layernorm vector width 1$\rightarrow$4, uncoalesced 64$\rightarrow$16; gemm/softmax no-transform. \\
\bottomrule
\end{tabularx}
\caption{Dashboard sections and their provenance modes. The mode is a
convention recorded by the aggregator; the report schema does not require it.}
\label{tab:dashboard}
\end{table}

\subsection{What the provenance does and does not enforce}
\label{sec:provenance}

The report schema (\Cref{app:report-schema}) requires \kst{kernel},
\kst{metrics}, \kst{limiter}, \kst{warnings}, and \kst{suggestions}. It does
not require \kst{schema_version} or \kst{provenance}; those are report-engine
and dashboard conventions. The per-section provenance object has a
\kst{mode} in \{measured, projected, no\_data\} and a boolean
\kst{measured}. The dashboard validator asserts the six sections, the boolean,
the evidence files, and the rendered module markers. The validator is not
called by \kst{ci.sh}; it is a separate check. The report states this because
a provenance convention that is described as enforced when it is only
intended is exactly the kind of claim this project is trying to avoid.

% ===========================================================================
